\documentclass[10pt]{extarticle}
\usepackage[a4paper,margin=1.5cm]{geometry}
\usepackage[T1]{fontenc}
\usepackage{amsmath,amssymb}
\usepackage{graphicx}
\usepackage[export]{adjustbox}
\usepackage{multicol}
\usepackage{enumitem}
\usepackage{booktabs,array,tabularx}
\setlength{\columnsep}{0.7cm}
\setlength{\parindent}{0pt}
\setlength{\parskip}{3pt}
\pagestyle{empty}
\makeatletter
% figure registry (figdefs.tex is generated by the extract step)
\newcommand{\qfigdef}[4]{\@namedef{qf@file@#1}{#2}\@namedef{qf@w@#1}{#3}\@namedef{qf@wide@#1}{#4}}
\newcommand{\qf@check}[1]{\@ifundefined{qf@file@#1}{\PackageError{score2}{Unknown figure id #1}{}}{}}
\newcommand{\qf@img}[4][]{\includegraphics[width=#2\dimexpr\@nameuse{qf@w@#4}pt\relax,max width=#3,max height=0.45\textheight,#1]{\@nameuse{qf@file@#4}}}
% Full-width questions: finalize wraps every question that contains a wide figure in \qfullbegin...\qfullend,
% so the WHOLE question (stem, figures, options) is set across the page in reading order, then the two columns
% resume. The question's head - its number, stem and first figure (\qfullsplit marks the end of the head) -
% is kept on one page, so a number is never left at a page foot with its figure overleaf; the rest may break
% across pages normally, so no big gaps are left. Without a \qfullsplit, the whole question is kept together
% if it is under 45% of a page. Never write these yourself.
\newif\ifqf@full
\newif\ifqf@boxed
\newsavebox{\qf@box}
\newcommand{\qf@open}{\global\qf@fulltrue\par\global\qf@boxedtrue
  \setbox\qf@box\vbox\bgroup\hsize\textwidth\linewidth\textwidth}
\newcommand{\qf@emit}[1]{\par\egroup\global\qf@boxedfalse
  \ifdim\dimexpr\ht\qf@box+\dp\qf@box\relax<#1\textheight\noindent\usebox{\qf@box}\par\else\unvbox\qf@box\fi}
\newcommand{\qfullsplit}{\ifqf@boxed\qf@emit{0.85}\fi}
\newcommand{\qf@close}{\ifqf@boxed\qf@emit{0.45}\fi\global\qf@fullfalse}
\newcommand{\qfullbegin}{\end{multicols}\qf@open}
\newcommand{\qfullend}{\qf@close\begin{multicols}{2}}
\newcommand{\qfullstart}{\qf@open}
\newcommand{\qfulldone}{\qf@close}
% \qfig[scale]{id}: block figure, centred, directly where it is written (size set by the tool).
% Fallback if a wide figure is not inside a full-width question: break out of the columns in place - no floats.
\newcommand{\qfig}[2][1]{\qf@check{#2}%
  \ifqf@full\qf@blk{#1}{#2}\else
  \ifnum\@nameuse{qf@wide@#2}=1 \expandafter\@firstoftwo\else\expandafter\@secondoftwo\fi
  {\end{multicols}\begin{center}\qf@img{#1}{\textwidth}{#2}\end{center}\begin{multicols}{2}}%
  {\qf@blk{#1}{#2}}\fi}
\newcommand{\qf@blk}[2]{\par\smallskip{\centering\qf@img{#1}{\linewidth}{#2}\par}\smallskip}
% \qfigi[scale]{id}: inline figure (inside option lists / tables / minipages); never breaks columns.
\newcommand{\qfigi}[2][1]{\qf@check{#2}\qf@img[valign=c]{#1}{\linewidth}{#2}}
\makeatother
% structure macros - lint reads these to check completeness and figure placement
\newcommand{\papersec}[1]{\par\bigskip{\centering\Large\bfseries PAPER-#1\par}\medskip}
\newcommand{\partsec}[2]{\par\medskip\noindent{\large\bfseries PART-#1: #2}\par\smallskip}
\newcommand{\sect}[1]{\par\smallskip\noindent\textbf{#1}\par}
\newcommand{\qn}[1]{\par\smallskip\noindent\textbf{#1.}~}
\newcommand{\paraq}[2]{\par\smallskip\noindent\textbf{Paragraph for Questions #1 and #2}\par}
\newenvironment{qoptions}{\begin{itemize}[leftmargin=1.8em,itemsep=1pt,topsep=2pt,label={}]}{\end{itemize}}
\newcommand{\opt}[2]{\item[(#1)] #2}
% figure registry (generated by score2_tools.py finalize)
\qfigdef{p1_phy_s1_q04_1}{Score Advanced - Test 05 (15 Jul 2021)-img/p1_phy_s1_q04_1.png}{100.9}{0}
\qfigdef{p1_phy_s1_q05_1}{Score Advanced - Test 05 (15 Jul 2021)-img/p1_phy_s1_q05_1.png}{124.6}{0}
\qfigdef{p1_phy_s1_q06_A}{Score Advanced - Test 05 (15 Jul 2021)-img/p1_phy_s1_q06_A.png}{91.1}{0}
\qfigdef{p1_phy_s1_q06_B}{Score Advanced - Test 05 (15 Jul 2021)-img/p1_phy_s1_q06_B.png}{77.4}{0}
\qfigdef{p1_phy_s1_q06_C}{Score Advanced - Test 05 (15 Jul 2021)-img/p1_phy_s1_q06_C.png}{75.7}{0}
\qfigdef{p1_phy_s1_q07_1}{Score Advanced - Test 05 (15 Jul 2021)-img/p1_phy_s1_q07_1.png}{122.7}{0}
\qfigdef{p1_phy_s1_q08_1}{Score Advanced - Test 05 (15 Jul 2021)-img/p1_phy_s1_q08_1.png}{126.4}{0}
\qfigdef{p1_phy_s1_q10_1}{Score Advanced - Test 05 (15 Jul 2021)-img/p1_phy_s1_q10_1.png}{123.5}{0}
\qfigdef{p1_phy_s1_q11_1}{Score Advanced - Test 05 (15 Jul 2021)-img/p1_phy_s1_q11_1.png}{124.7}{0}
\qfigdef{p1_phy_s1_q12_A}{Score Advanced - Test 05 (15 Jul 2021)-img/p1_phy_s1_q12_A.png}{38.5}{0}
\qfigdef{p1_phy_s1_q12_B}{Score Advanced - Test 05 (15 Jul 2021)-img/p1_phy_s1_q12_B.png}{37.3}{0}
\qfigdef{p1_phy_s1_q12_C}{Score Advanced - Test 05 (15 Jul 2021)-img/p1_phy_s1_q12_C.png}{38.0}{0}
\qfigdef{p1_phy_s1_q12_D}{Score Advanced - Test 05 (15 Jul 2021)-img/p1_phy_s1_q12_D.png}{40.2}{0}
\qfigdef{p1_phy_s1_q14_1}{Score Advanced - Test 05 (15 Jul 2021)-img/p1_phy_s1_q14_1.png}{126.3}{0}
\qfigdef{p1_phy_s1_q15_1}{Score Advanced - Test 05 (15 Jul 2021)-img/p1_phy_s1_q15_1.png}{76.9}{0}
\qfigdef{p1_phy_s1_q16_1}{Score Advanced - Test 05 (15 Jul 2021)-img/p1_phy_s1_q16_1.png}{95.5}{0}
\qfigdef{p1_phy_s1_q17_1}{Score Advanced - Test 05 (15 Jul 2021)-img/p1_phy_s1_q17_1.png}{101.8}{0}
\qfigdef{p1_phy_s1_q20_1}{Score Advanced - Test 05 (15 Jul 2021)-img/p1_phy_s1_q20_1.png}{96.7}{0}
\qfigdef{p1_phy_s2_q02_1}{Score Advanced - Test 05 (15 Jul 2021)-img/p1_phy_s2_q02_1.png}{134.4}{0}
\qfigdef{p1_phy_s2_q05_A}{Score Advanced - Test 05 (15 Jul 2021)-img/p1_phy_s2_q05_A.png}{105.5}{0}
\qfigdef{p1_phy_s2_q05_B}{Score Advanced - Test 05 (15 Jul 2021)-img/p1_phy_s2_q05_B.png}{104.9}{0}
\qfigdef{p1_phy_s2_q05_C}{Score Advanced - Test 05 (15 Jul 2021)-img/p1_phy_s2_q05_C.png}{103.6}{0}
\qfigdef{p1_phy_s2_q06_1}{Score Advanced - Test 05 (15 Jul 2021)-img/p1_phy_s2_q06_1.png}{97.2}{0}
\qfigdef{p1_phy_s2_q09_1}{Score Advanced - Test 05 (15 Jul 2021)-img/p1_phy_s2_q09_1.png}{78.7}{0}
\qfigdef{p1_phy_s2_q10_1}{Score Advanced - Test 05 (15 Jul 2021)-img/p1_phy_s2_q10_1.png}{97.2}{0}
\qfigdef{p1_chem_s1_q02_A}{Score Advanced - Test 05 (15 Jul 2021)-img/p1_chem_s1_q02_A.png}{44.6}{0}
\qfigdef{p1_chem_s1_q02_B}{Score Advanced - Test 05 (15 Jul 2021)-img/p1_chem_s1_q02_B.png}{45.5}{0}
\qfigdef{p1_chem_s1_q02_C}{Score Advanced - Test 05 (15 Jul 2021)-img/p1_chem_s1_q02_C.png}{51.5}{0}
\qfigdef{p1_chem_s1_q02_D}{Score Advanced - Test 05 (15 Jul 2021)-img/p1_chem_s1_q02_D.png}{54.1}{0}
\qfigdef{p1_chem_s1_q05_1}{Score Advanced - Test 05 (15 Jul 2021)-img/p1_chem_s1_q05_1.png}{115.1}{0}
\qfigdef{p1_chem_s1_q08_1}{Score Advanced - Test 05 (15 Jul 2021)-img/p1_chem_s1_q08_1.png}{105.8}{0}
\qfigdef{p1_chem_s1_q14_A}{Score Advanced - Test 05 (15 Jul 2021)-img/p1_chem_s1_q14_A.png}{52.7}{0}
\qfigdef{p1_chem_s1_q14_B}{Score Advanced - Test 05 (15 Jul 2021)-img/p1_chem_s1_q14_B.png}{58.5}{0}
\qfigdef{p1_chem_s1_q14_C}{Score Advanced - Test 05 (15 Jul 2021)-img/p1_chem_s1_q14_C.png}{58.5}{0}
\qfigdef{p1_chem_s1_q14_D}{Score Advanced - Test 05 (15 Jul 2021)-img/p1_chem_s1_q14_D.png}{56.5}{0}
\qfigdef{p1_chem_s1_q19_1}{Score Advanced - Test 05 (15 Jul 2021)-img/p1_chem_s1_q19_1.png}{150.6}{0}
\qfigdef{p1_chem_s1_q19_B}{Score Advanced - Test 05 (15 Jul 2021)-img/p1_chem_s1_q19_B.png}{43.1}{0}
\qfigdef{p1_chem_s1_q20_1}{Score Advanced - Test 05 (15 Jul 2021)-img/p1_chem_s1_q20_1.png}{75.7}{0}
\raggedcolumns
\makeatletter\@beginparpenalty=10000\makeatother
\newcolumntype{L}[1]{>{\raggedright\arraybackslash}p{#1}}
\begin{document}
\begin{center}{\Large\bfseries Score Advanced --- Test 05 (15 Jul 2021)}\end{center}
\begin{multicols}{2}
\raggedcolumns
\papersec{1}
\partsec{1}{PHYSICS}
\sect{SEC-1}
\qn{1} A shell of mass m is at rest initially. It explodes into three fragements having masses in the ratio 2 : 2 : 1. The fragements having equal masses fly off along mutually perpendicular directions with speed v. What will be the speed of the ligher fragments?
\begin{qoptions}
\opt{A}{v}
\opt{B}{$\sqrt{2}$v}
\opt{C}{$2\sqrt{2}$v}
\opt{D}{$3\sqrt{2}$v}
\end{qoptions}
\qn{2} In a P--N Junction diode not connected to any circuit
\begin{qoptions}
\opt{A}{potential is the same every where.}
\opt{B}{the P - type side is at a higher potential than the N - type side.}
\opt{C}{there is an electric field at the junction directed from the N - type side to the P-type side.}
\opt{D}{there is an electric field at the junction directed from the P - type side to the N-type side.}
\end{qoptions}
\qn{3} Mark the correct statement
\begin{qoptions}
\opt{A}{The energy contained in a small volume through which an electromagnetic wave is passing oscillates with the frequency of the wave}
\opt{B}{Magnetic energy has zero average value in a plane electromagnetic wave}
\opt{C}{Intensity in electromagnetic wave is given by $\dfrac{1}{2}\varepsilon_0 E_0^2 c$ where E$_0$ is maximum electric field}
\opt{D}{Energy density of electromagnetic wave is given by $\dfrac{1}{4}\varepsilon_0 E_0^2$}
\end{qoptions}
\qn{4} Find the equivalent capacitance of circuit and charge on 5$\mu$F capacitor :\par
\qfig{p1_phy_s1_q04_1}
\begin{qoptions}
\opt{A}{4$\mu$F, 50$\mu$C}
\opt{B}{8$\mu$F, 25$\mu$C}
\opt{C}{4$\mu$F, 25$\mu$C}
\opt{D}{8$\mu$F, 50$\mu$C}
\end{qoptions}
\qn{5} An electric dipole is projected with velocity v in a uniform magnetic field directed inside the plane of paper as shown. Choose the correct option.\par
\qfig{p1_phy_s1_q05_1}
\begin{qoptions}
\opt{A}{Dipole will experience a net force in case-I due to external magnetic field}
\opt{B}{Dipole will experience a net force in case-II due to external magnetic field}
\opt{C}{Dipole will experience a net torque in case-I due to external magnetic field}
\opt{D}{Dipole will experience a net torque in case-II due to external magnetic field}
\end{qoptions}
\qn{6} The time variations of signals are given as in A, B and C. Point out the true statement from the following.\par
\qfig{p1_phy_s1_q06_A}
\qfig{p1_phy_s1_q06_B}
\qfig{p1_phy_s1_q06_C}
\begin{qoptions}
\opt{A}{A, B and C are analogue signals}
\opt{B}{A and C digital, but B is analogue signal}
\opt{C}{A and C are analogue but B is digital signal}
\opt{D}{A, B and C are digital signal}
\end{qoptions}
\qn{7} If the P-V diagrams of two thermodynamics devices working in a cyclic process are as shown in the figure, then :-\par
\qfig{p1_phy_s1_q07_1}
\begin{qoptions}
\opt{A}{A is a heat engine, B is a heat pump/refrigerator}
\opt{B}{B is a heat engine, A is a heat pump/refrigerator}
\opt{C}{both A and B are heat engines}
\opt{D}{both A and B are heat pumps/refrigerator}
\end{qoptions}
\qn{8} Figure shows a colour --coded resistor. What is the resistance of this resistor?\par
\qfig{p1_phy_s1_q08_1}
\begin{qoptions}
\opt{A}{47 $\times$ 10$^1$ $\Omega$ $\pm$5\%}
\opt{B}{23 $\times$ 10$^1$ $\Omega$ $\pm$5\%}
\opt{C}{47 $\times$ 10$^2$ $\Omega$ $\pm$5\%}
\opt{D}{23 $\times$ 10$^2$ $\Omega$ $\pm$10\%}
\end{qoptions}
\qn{9} Select the correct statement :-
\begin{qoptions}
\opt{A}{Total internal reflection takes place when light travel from rarer to denser medium.}
\opt{B}{Focal length of a convex lens is positive irrespective of medium in which it is placed.}
\opt{C}{Focal length of spherical mirror depends on its radius of curvature not on the surrounding medium.}
\opt{D}{Blue colour of sky appear due to total internal reflection}
\end{qoptions}
\qn{10} A glass hemisphere of radius R and material having refractive index 1.5 is silvered on its flat face as shown in the figure. A small object of height h is located at a distance 2R from surface of hemisphere. The final image will form :-\par
\qfig{p1_phy_s1_q10_1}
\begin{qoptions}
\opt{A}{at a distance of R from silvered surface, on right side}
\opt{B}{on the hemispherical surface}
\opt{C}{at the plane surface}
\opt{D}{at a distance of 2 R from the curved surface on left side}
\end{qoptions}
\qn{11} A solid sphere of mass M and radius r slips on a rough horizontal plane. At some instant it has translational velocity v$_0$ and rotational velocity about the centre v$_0$/2r. The translational velocity after the sphere starts pure rolling :\par
\qfig{p1_phy_s1_q11_1}
\begin{qoptions}
\opt{A}{6v$_0$/ 7 in forward direction}
\opt{B}{6v$_0$/ 7 in backward direction}
\opt{C}{7v$_0$/ 6 in forward direction}
\opt{D}{6v$_0$/ 6 in backward direction}
\end{qoptions}
\qn{12} Which of the following figures best represents the stable equilibrium of a uniform semicircular rod in vertical plane. It is hinged at point A.
\begin{qoptions}
\opt{A}{\qfigi{p1_phy_s1_q12_A}}
\opt{B}{\qfigi{p1_phy_s1_q12_B}}
\opt{C}{\qfigi{p1_phy_s1_q12_C}}
\opt{D}{\qfigi{p1_phy_s1_q12_D}}
\end{qoptions}
\qn{13} Two radioactive sources A and B of half lives of 2hr and 4hr respectively, initially contain the same number of radioactive atoms. At the end of 2 hours, their rates of disintegration are in the ratio :-
\begin{qoptions}
\opt{A}{4 : 1}
\opt{B}{2 : 1}
\opt{C}{$\sqrt{2}$ : 1}
\opt{D}{1 : 1}
\end{qoptions}
\qn{14} One end of a thermally insulated rod is kept at a temperature T$_1$ and the other at T$_2$. The rod is composed of two sections of lengths $\ell_1$ and $\ell_2$ and thermal conductivities K$_1$ and K$_2$ respectively. The temperature at the interface of the two sections is-\par
\qfig{p1_phy_s1_q14_1}
\begin{qoptions}
\opt{A}{(K$_2\ell_2$T$_1$ + K$_1\ell_1$T$_2$)/(K$_1\ell_2$ + K$_2\ell_1$)}
\opt{B}{(K$_2\ell_1$T$_1$ + K$_1\ell_2$T$_2$)/(K$_2\ell_1$ + K$_1\ell_2$)}
\opt{C}{(K$_1\ell_2$T$_1$ + K$_2\ell_1$T$_2$)/(K$_1\ell_2$ + K$_2\ell_1$)}
\opt{D}{(K$_1\ell_1$T$_1$ + K$_2\ell_2$T$_2$)/(K$_1\ell_1$ + K$_2\ell_2$)}
\end{qoptions}
\qn{15} A conducting sphere of radius R and a concentric thick spherical shell of inner radius 2R and outer radius 3R is shown in figure. A charge +10Q is given to the shell and inner sphere is earthed. Then charge on inner sphere is :\par
\qfig{p1_phy_s1_q15_1}
\begin{qoptions}
\opt{A}{--4Q}
\opt{B}{--10Q}
\opt{C}{zero}
\opt{D}{none}
\end{qoptions}
\qn{16} Two vessels A and B of different shapes have the same base area and are filled with water up to the same height h (see figure). The force exerted by water on the base is F$_A$ for vessel A and F$_B$ for vessel B. The respective weights of the water filled in vessels are W$_A$ and W$_B$. Then\par
\qfig{p1_phy_s1_q16_1}
\begin{qoptions}
\opt{A}{F$_A$ $>$ F$_B$ ; W$_A$ $>$ W$_B$}
\opt{B}{F$_A$ = F$_B$ ; W$_A$ $>$ W$_B$}
\opt{C}{F$_A$ = F$_B$ ; W$_A$ $<$ W$_B$}
\opt{D}{F$_A$ $>$ F$_B$ ; W$_A$ = W$_B$}
\end{qoptions}
\qn{17} The graph shows current v/s voltage in a series RLC A.C. circuit. The arrow indicates the direction that this curve is drawn as time progresses. In this plot, the\par
\qfig{p1_phy_s1_q17_1}
\begin{qoptions}
\opt{A}{current lags the voltage by about 90 degrees}
\opt{B}{current leads the voltage by about 90 degrees}
\opt{C}{current and voltage are in phase}
\opt{D}{current and voltage are 180 degrees out of phase}
\end{qoptions}
\qn{18} When monochromatic light falls normally on the surface of the soap film in air, the intensity of the reflected light depends on the wavelength: it has a maximum at $\lambda_1$ = 630nm and the minimum at $\lambda_2$ = 525 nm. What can be the film thickness? Refractive index of the film n = 4/3.
\begin{qoptions}
\opt{A}{590.625 nm}
\opt{B}{640 nm}
\opt{C}{562.625 nm}
\opt{D}{466.66 nm}
\end{qoptions}
\qn{19} In a resonance tube experiment, an 80 cm air column is in resonance with a tuning fork in first overtone. Which equation can represent correct pressure variation in the air column (x = 0 is the top point of the tube, neglect end correction, speed of sound = 320 m/sec)
\begin{qoptions}
\opt{A}{A sin$\dfrac{15\pi}{8}$x cos 600$\pi$t}
\opt{B}{A cos$\dfrac{15\pi}{8}$x sin 600 $\pi$t}
\opt{C}{A cos$\dfrac{15\pi}{8}$x sin 300 $\pi$t}
\opt{D}{A sin$\dfrac{15\pi}{8}$x sin 300 $\pi$t}
\end{qoptions}
\qn{20} The variation of density ($\rho$) with temperature (T) of an ideal gas of molecular mass (M) at constant pressure is plotted as shown in figure. The pressure of the gas at state P will be :-\par
\qfig{p1_phy_s1_q20_1}
\begin{qoptions}
\opt{A}{$\dfrac{\text{RT}^2}{\sqrt{3}\text{M}}$}
\opt{B}{$\dfrac{\sqrt{3}\text{RT}^2}{\text{M}}$}
\opt{C}{$\dfrac{\text{RT}^2}{\text{M}}$}
\opt{D}{$\dfrac{\text{RT}^2}{2\text{M}}$}
\end{qoptions}
\sect{SEC-2}
\qn{1} The acceleration due to gravity (g) at a place is determined by using a simple pendulum of length $\ell$ = (100 $\pm$ 0.1) cm. If the time period is T = (2 $\pm$ 0.01) s. Then maximum percentage error in the measurement of g.
\qn{2} Fill the value of reading of spring balance (in N) as shown in figure. When the wedge of mass 2kg is in equilibrium and all surfaces are smooth.\par
\qfig{p1_phy_s2_q02_1}
\qn{3} 1000 small mercury (Hg) drops each of the radius r and charge q coalesce together to form one spherical drop. The potential of the bigger drop is 10$^n$ times larger than of the smaller ones then the value of n is
\qn{4} A particle moves along a curve y = $\dfrac{x^2}{2}$ and x = $\dfrac{t^2}{2}$, where x and y are measured in meter and t is in seconds. Then the speed of particle at t = 2 sec is v m/s then $\dfrac{v^2}{80}$ is
\qn{5} In the following combination of logic gates, the outputs of A, B and C are x, y and z respectively then the value of x + y + z is\par
\qfig{p1_phy_s2_q05_A}
\qfig{p1_phy_s2_q05_B}
\qfig{p1_phy_s2_q05_C}
\qn{6} In the circuit show below, the key K is closed at t = 0. The current through the battery at t = 0 is xA and at t = $\infty$ is yA, then x + y is\par
\qfig{p1_phy_s2_q06_1}
\qn{7} A tunnel is dug inside the earth across one of its diameters. A particle is projected inside the tunnel with velocity $\sqrt{\dfrac{\text{GM}}{\text{R}}}$ from one of its ends. Then maximum velocity attained by the particle in the subsequent motion is $\sqrt{\dfrac{\text{nGM}}{\text{R}}}$, then value of n. (assuming tunnel to be frictionless)
\qn{8} If light of wavelength of maximum intensity emitted from surface at temperature T$_1$ is used to cause photoelectric emission from a metallic surface, the maximum kinetic energy of the emitted electron is 6eV, which is 3 times the work function of the metallic surface. If light of wavelength of maximum intensity emitted from a surface at temperature T$_2$ (T$_2$ = 2T$_1$) is used, the maximum kinetic energy of the photoelectrons emitted (in eV) is :-
\qn{9} A spring-block system is arranged as shown in diagram. Its time period in terms of mass of block `M' and spring of spring constant `K' is $2\pi\sqrt{\dfrac{\text{nM}}{\text{K}}}$, then value of n. (Assume all pulleys to be ideal)\par
\qfig{p1_phy_s2_q09_1}
\qn{10} Three infinite wires are arranged in space in three dimensions (along x, y and z axis) as shown. Each wire carries current i. Then magnetic field at A is $\dfrac{\text{m}\mu_0 i}{\pi r}\left(\hat{i}+\hat{j}-n\hat{k}\right)$ then the value of $\dfrac{\text{n}}{\text{m}}$ is\par
\qfig{p1_phy_s2_q10_1}
\partsec{2}{CHEMISTRY}
\sect{SEC-1}
\qn{1} In which of the following options order of arrangement does not agree with the variation in property indicated against it.
\begin{qoptions}
\opt{A}{Al$^{+3}$ $<$ Mg$^{+2}$ $<$ Na$^+$ $<$ F$^-$ (Increasing ionic size)}
\opt{B}{B$^+$ $<$ C$^+$ $<$ N$^+$ $<$ O$^+$ (Ionisation enthalpy)}
\opt{C}{I $<$ Br $<$ F $<$ Cl (Increasing $|\Delta$Heg$|$)}
\opt{D}{Li $<$ Na $<$ K $<$ Rb (Increasing metallic radius)}
\end{qoptions}
\qn{2} Benzoic acid is heated with PCl$_5$ and then reduced by H$_2$ in presence of Pd/BaSO$_4$. A is formed which is heated with RNH$_2$, B is formed which is reduced by H$_2$/Ni catalyst, C is formed as the final product. C in the reaction sequence is :
\begin{qoptions}
\opt{A}{\qfigi{p1_chem_s1_q02_A}}
\opt{B}{\qfigi{p1_chem_s1_q02_B}}
\opt{C}{\qfigi{p1_chem_s1_q02_C}}
\opt{D}{\qfigi{p1_chem_s1_q02_D}}
\end{qoptions}
\qn{3} Which alkene on heating with alkaline KMnO$_4$ solution gives acetone and a gas, which turns lime water milky?
\begin{qoptions}
\opt{A}{2-methyl 2-butene}
\opt{B}{Isobutene}
\opt{C}{1-butene}
\opt{D}{2-butene}
\end{qoptions}
\qn{4} \par\noindent CaCO$_3$ $\xrightarrow{\Delta}$ CO$_2$ + X\\
X + coke $\xrightarrow{\text{High temperature}}$ Y\\
Y + H$_2$O $\longrightarrow$ Z\\
Z $\xrightarrow[\text{H}_3\text{O}^+]{\text{HgSO}_4}$ Z$_3$\\
Z $\xrightarrow[\text{Pt}]{\text{H}_2\ (1\text{ eq})}$ Z$_1$ $\xrightarrow{\text{HBr}}$ Z$_2$\\
Z$_1$, Z$_2$, Z$_3$ are respectively :
\begin{qoptions}
\opt{A}{C$_2$H$_4$, C$_2$H$_5$Br, CH$_3$CH$_2$CHO}
\opt{B}{C$_2$H$_4$, C$_2$H$_5$Br, CH$_3$CHO}
\opt{C}{C$_3$H$_6$, isopropyl bromide, CH$_3$COCH$_3$}
\opt{D}{C$_3$H$_6$, n-propyl bromide, CH$_3$COCH$_3$}
\end{qoptions}
\qn{5} \qfig{p1_chem_s1_q05_1}
Product is example of
\begin{qoptions}
\opt{A}{$\gamma$-Lactone}
\opt{B}{$\delta$-Lactone}
\opt{C}{Cyclic hemiacetal}
\opt{D}{$\alpha,\beta$-unsaturated acid}
\end{qoptions}
\qn{6} The incorrect match of the items given in column-I with examples given in column-II
\begin{qoptions}
\opt{A}{Hydrogen bond is present \quad (a) H$_2$}
\opt{B}{Resonance is involved \quad (b) CO$_2$}
\opt{C}{Ionic compound \quad (c) NaCl}
\opt{D}{Covalent solid \quad (d) C}
\end{qoptions}
\qn{7} The colour of the co-ordination compounds depends on the crystal field splitting. What will be the correct order of absorption of wavelength of light in the visible region.
\begin{qoptions}
\opt{A}{[Co(NH$_3$)$_6$]$^{+3}$ $>$ [Co(CN)$_6$]$^{-3}$ $>$ [Co(H$_2$O)$_6$]$^{+3}$}
\opt{B}{[Co(NH$_3$)$_6$]$^{+3}$ $>$ [Co(H$_2$O)$_6$]$^{+3}$ $>$ [Co(CN)$_6$]$^{-3}$}
\opt{C}{[Co(H$_2$O)$_6$]$^{+3}$ $>$ [Co(NH$_3$)$_6$]$^{+3}$ $>$ [Co(CN)$_6$]$^{-3}$}
\opt{D}{[Co(CN)$_6$]$^{-3}$ $>$ [Co(NH$_3$)$_6$]$^{+3}$ $>$ [Co(H$_2$O)$_6$]$^{+3}$}
\end{qoptions}
\qn{8} Energy level diagram of formation and combustion of components in standard state is given-\par
\qfig{p1_chem_s1_q08_1}
Magnitude of X + Y is equal to -
\begin{qoptions}
\opt{A}{$\Delta$H$^0_{\text{c}}$(C) + $\Delta$H$^0_{\text{c}}$(H$_2$)}
\opt{B}{$\Delta$H$^0_{\text{f}}$(CO$_2$) + $\Delta$H$^0_{\text{f}}$(H$_2$O)}
\opt{C}{$\Delta$H$^0_{\text{f}}$(glucose) + $\Delta$H$^0_{\text{c}}$(glucose)}
\opt{D}{$\Delta$H$^0_{\text{f}}$(glucose) + $\Delta$H$^0_{\text{f}}$(CO$_2$) + $\Delta$H$^0_{\text{f}}$(H$_2$O)}
\end{qoptions}
\qn{9} Match the coordination compounds given in column-I with the central metal atom given in column-II and assign the correct code.\par
\noindent{\footnotesize\setlength{\tabcolsep}{2pt}\begin{tabular}{|L{.4cm}|L{3cm}|L{.4cm}|L{2.4cm}|}
\hline
 & Column-I Co-ordination compounds & & Column-II (Central metal) \\ \hline
a & Chlorophyll & p & Rhodium \\ \hline
b & Blood pigment & q & Cobalt \\ \hline
c & Wilkinson's catalyst & r & Calcium \\ \hline
d & Vitamin B$_{12}$ & s & Iron \\ \hline
 & & t & Magnesium \\ \hline
\end{tabular}}\par
\begin{qoptions}
\opt{A}{a-(t);b-(s);c-(p);d-(q)}
\opt{B}{a-(r);b-(s);c-(t);d-(p)}
\opt{C}{a-(s);b-(r);c-(q);d-(p)}
\opt{D}{a-(t);b-(s);c-(r) ;d-(q)}
\end{qoptions}
\qn{10} Which of the following change is not possible during roasting.
\begin{qoptions}
\opt{A}{CaCO$_3$ $\longrightarrow$ CaO + CO$_2$}
\opt{B}{Al$_2$O$_3$.xH$_2$O $\longrightarrow$ Al$_2$O$_3$ + xH$_2$O}
\opt{C}{Fe$_2$O$_3$ + O$_2$ $\longrightarrow$ Fe$_3$O$_4$}
\opt{D}{ZnS + $\dfrac{3}{2}$O$_2$ $\longrightarrow$ ZnO + SO$_2$}
\end{qoptions}
\qn{11} Which of the following reagent can not be used to distinguish p-chloro aniline and anilinium hydrochloride?
\begin{qoptions}
\opt{A}{PhSO$_2$Cl + NaOH}
\opt{B}{AgNO$_3$}
\opt{C}{NaHCO$_3$}
\opt{D}{All}
\end{qoptions}
\qn{12} An organic compound A having C, H, O 0.108 gm A on combustion gives 0.308 gm CO$_2$ and 0.072 gm H$_2$O. A is insoluble in NaOH and NaHCO$_3$. A reacts with HI to give compound C and D. D reacts with neutral FeCl$_3$ to give violet color\par
\noindent D $\xrightarrow[\text{2--4 atm/H}^+]{\text{NaOH+CO}_2}$ D$_1$ $\xrightarrow{\text{CH}_3\text{COCl}}$ D$_2$\par
The product D$_2$ is
\begin{qoptions}
\opt{A}{An antiallergic}
\opt{B}{An antiseptic}
\opt{C}{An antibiotic}
\opt{D}{An antipyretic}
\end{qoptions}
\qn{13} Which of the following reactions increase production of dihydrogen from synthesis gas?
\begin{qoptions}
\opt{A}{CH$_4$(s) + H$_2$O(g) $\xrightarrow[\text{Ni}]{1270\text{ K}}$ CO + 3H$_2$(g)}
\opt{B}{C(s) + H$_2$O(g) $\xrightarrow{1270\text{ K}}$ CO(g) + H$_2$(g)}
\opt{C}{CO + H$_2$O(g) $\xrightarrow[\text{Catalyst}]{673\text{ K}}$ CO$_2$(g) + H$_2$(g)}
\opt{D}{C$_2$H$_6$ + 2H$_2$O $\xrightarrow[\text{Ni}]{1270\text{ K}}$ 2CO(g) + 5H$_2$(g)}
\end{qoptions}
\qn{14} X is a hexose. On catalytic reduction X gives two hexahydric alcohol Y and Z. Y can be obtained from D(+) glucose by reducing with NaBH$_4$. X can be :
\begin{qoptions}
\opt{A}{\qfigi{p1_chem_s1_q14_A}}
\opt{B}{\qfigi{p1_chem_s1_q14_B}}
\opt{C}{\qfigi{p1_chem_s1_q14_C}}
\opt{D}{\qfigi{p1_chem_s1_q14_D}}
\end{qoptions}
\qn{15} Terylene (Dacron) is a condensation polymer of :
\begin{qoptions}
\opt{A}{Formaldehyde and urea}
\opt{B}{Ethylene glycol and ethylene diisocyanate}
\opt{C}{Ethylene glycol and phthalic acid}
\opt{D}{Ethylene glycol and terephthalic acidl}
\end{qoptions}
\qn{16} For acid catalysed hydrolysis of ester, rate constant is given as K= $\dfrac{1}{t}\ln\left(\dfrac{V_\infty - V_0}{V_\infty - V_t}\right)$ where V$_0$, V$_t$ and V$_\infty$ are volume of standard NaOH solution required to neutralise acid present at given time. If ester gets 50\% hydrolysed then -
\begin{qoptions}
\opt{A}{V$_\infty$ = V$_t$}
\opt{B}{V$_\infty$ = V$_t$--V$_0$}
\opt{C}{V$_\infty$ = 2V$_t$--V$_0$}
\opt{D}{V$_\infty$ = V$_t$ --2V$_0$}
\end{qoptions}
\qn{17} Silicon derivatives has a strong tendency to form polymeric resins like silicones. The chain length of silicone polymer can be controlled or stopped by adding.
\begin{qoptions}
\opt{A}{MeSiCl$_3$}
\opt{B}{Me$_2$SiCl$_2$}
\opt{C}{Me$_3$SiCl}
\opt{D}{Me$_4$Si}
\end{qoptions}
\qn{18} Air pollutants that produce photochemical smog are :
\begin{qoptions}
\opt{A}{Ozone, chloride and SO$_2$}
\opt{B}{O$_2$, Cl$_2$, HNO$_3$}
\opt{C}{Nitrogen oxide, PAN and Acrolein}
\opt{D}{CO$_2$, CO and SO$_2$}
\end{qoptions}
\qn{19} \qfig{p1_chem_s1_q19_1}
Incorrect statement regarding above compounds is
\begin{qoptions}
\opt{A}{Basic strength I $>$ II $>$ III}
\opt{B}{Compound III is more basic than \qfigi{p1_chem_s1_q19_B}}
\opt{C}{All are examples of heterocyclic aromatic compounds}
\opt{D}{Compound II shows inductive, hyperconjugation, resonance effect.}
\end{qoptions}
\qn{20} \qfig{p1_chem_s1_q20_1}
Identify the best reagent sequence for above conversion :
\begin{qoptions}
\opt{A}{Cl$_2$ / h$\nu$; alcoholic KOH; O$_3$ / Zn; H$_2$O}
\opt{B}{Cl$_2$ / h$\nu$; tertiarybutoxide ion; O$_3$/Zn, H$_2$O}
\opt{C}{NBS, alcoholic KOH, $\Delta$, acidic KMnO$_4$}
\opt{D}{SOCl$_2$, tertiary butoxide ion, alkaline KMnO$_4$}
\end{qoptions}
\sect{SEC-2}
\qn{1} Volume of water (in ml) added in 2ml, 63\% w/w concentrated acidic HNO$_3$ solution of density 1.5 gm/ml to make its final solution 1N is -
\qn{2} Moles of K$_2$Cr$_2$O$_7$ required to oxidise 1 mole of Cu$_3$P to form CuSO$_4$ and H$_3$PO$_4$ in acidic medium is-
\qn{3} How many of the following specie(s) is/are both non linear and non-polar:-\par
H$_2$O, HOCl, BeCl$_2$, Cl$_2$O, I$_3^{-}$, N$_3^{-}$, IF$_2^{+}$
\qn{4} The radius of two of the first four Bohr's orbits in hydrogen atom are in ratio 1 : 4. If the energy difference (in eV) between them is x \& y. Give (x + y).
\qn{5} The ratio of Boyle's temperature and critical temperature of a real gas is -
\qn{6} Two moles of ideal gas (C$_{\text{v , m}}$ = $\dfrac{3}{2}$R) is heated from 300K to 600K at constant pressure. If the change in entropy of gas ($\Delta$S) is xRln2. Give the value of `x'.
\qn{7} Consider following gaseous reactions at equilibrium :\par
2AB $\rightleftharpoons$ A$_2$ + B$_2$, \quad K = 2.5 $\times$ 10$^5$\par
AB + $\dfrac{1}{2}$C$_2$ $\rightleftharpoons$ ABC, K = 5 $\times$ 10$^{-3}$\par
If the value of equilibrium constant of following reaction is x $\times$ 10$^y$. Give (x + y) :\par
$\dfrac{1}{2}$A$_2$ + $\dfrac{1}{2}$B$_2$ + $\dfrac{1}{2}$C$_2$ $\rightleftharpoons$ ABC
\qn{8} If the solubility of Al(OH)$_3$, (K$_{\text{sp}}$ = 10$^{-33}$) in a buffer solution of pH = 4 is 10$^x$. Give x?
\qn{9} Metal carbonates decompose on heating to give metal oxide and carbon dioxide. Which of the following metal carbonates is most stable thermally?\par
(i) MgCO$_3$ \quad (ii) CaCO$_3$ \quad (iii) SrCO$_3$ \quad (iv) BaCO$_3$\par
Represent your answer as the period of the metal involved in it.
\qn{10} The time in hours taken to decompose electrolytically 18 gm of water by 2 ampere current is about -
\partsec{3}{MATHEMATICS}
\sect{SEC-1}
\qn{1} Let tangents are drawn from point (x,6), (where x $<$ 0), to parabola y$^2$ = 8x. If chord of contact subtends right angle at vertex of parabola then mid point of chord of contact is
\begin{qoptions}
\opt{A}{(6,17)}
\opt{B}{(17,6)}
\opt{C}{(8,12)}
\opt{D}{(5,12)}
\end{qoptions}
\qn{2} If 64sin$^6\alpha$ + 729 cos$^6\beta$ + 4 = 36 sin$\alpha$ cos$\beta$ where $\alpha,\beta \in \left(0,\dfrac{\pi}{2}\right)$ then value of $\dfrac{\tan\beta}{\tan\alpha}$ is
\begin{qoptions}
\opt{A}{$18\sqrt{3}$}
\opt{B}{34}
\opt{C}{$13\sqrt{8}$}
\opt{D}{24}
\end{qoptions}
\qn{3} If p $\rightarrow$ ($\sim$ p$\vee\sim$ q) is false then truth values of p and q are respectively
\begin{qoptions}
\opt{A}{F, F}
\opt{B}{F, T}
\opt{C}{T, T}
\opt{D}{T, F}
\end{qoptions}
\qn{4} If f(x) = $\sin^{-1}\left(\dfrac{2.3^x}{1+9^x}\right)$, then f$'\left(-\dfrac{1}{2}\right)$ equals
\begin{qoptions}
\opt{A}{$\sqrt{3}\log_e\sqrt{3}$}
\opt{B}{$-\sqrt{3}\log_e 3$}
\opt{C}{$-\sqrt{3}\log_e\sqrt{3}$}
\opt{D}{$\sqrt{3}\log_e 3$}
\end{qoptions}
\qn{5} Let f(x) = $\dfrac{(x-2)}{4}+\dfrac{(x-2)^3}{12}+\dfrac{(x-2)^5}{20}+\dfrac{(x-2)^7}{28}+\ldots$ upto $\infty$ terms where 1 $<$ x $<$ 3, then f$'$(x) is
\begin{qoptions}
\opt{A}{$\dfrac{1}{4(x-1)(3-x)}$}
\opt{B}{$\dfrac{1}{4(x-2)^2}$}
\opt{C}{$\dfrac{1}{4(3-x)}$}
\opt{D}{$\dfrac{1}{(x-3)(x-1)}$}
\end{qoptions}
\qn{6} If x $\in$ [0, 2$\pi$] and $\begin{vmatrix} 0 & \cos x & -\sin x\\ \sin x & 0 & \cos x\\ \cos x & \sin x & 0\end{vmatrix} = 0$, then $\sum\tan\left(\dfrac{\pi}{3}+x\right)$ is equal to-
\begin{qoptions}
\opt{A}{$\sqrt{3}-2$}
\opt{B}{$4+2\sqrt{3}$}
\opt{C}{$-2(2+\sqrt{3})$}
\opt{D}{$4-2\sqrt{3}$}
\end{qoptions}
\qn{7} If $\dfrac{dy}{dx} = \dfrac{y^3}{e^{2x}+y^2}$ and y(0) = 1, then solution of differential equation is-
\begin{qoptions}
\opt{A}{y$^2$ = e$^{2x}$ -- 2e$^{2x}\ln$y}
\opt{B}{y$^2$ = e$^{2x}$ -- $\dfrac{1}{2}$e$^{2x}$. $\ln$y}
\opt{C}{y$^2$ = e$^{2x}$ + 2e$^{2x}\ln$y}
\opt{D}{y$^2$ = e$^{2x}$ + $\dfrac{1}{2}$e$^{2x}\ln$y}
\end{qoptions}
\qn{8} Let $\displaystyle\int_0^1 \tan^{-1}\left(\dfrac{\tan x}{2}\right)dx = k$, then value of $\displaystyle\int_0^1 \tan^{-1}\left(\dfrac{\tan^2 x - 2}{3\tan x}\right)dx$ is
\begin{qoptions}
\opt{A}{$\pi - k + \dfrac{1}{2}$}
\opt{B}{$k - \dfrac{\pi}{2} - 1$}
\opt{C}{$k - \pi + 1$}
\opt{D}{$k - \dfrac{\pi}{2} + \dfrac{1}{2}$}
\end{qoptions}
\qn{9} Shortest distance between y$^2$ = x$^3$ and curve 9x$^2$ + 9y$^2$ -- 30y + 16 = 0
\begin{qoptions}
\opt{A}{$\dfrac{\sqrt{13}}{3} - 1$}
\opt{B}{$\dfrac{2\sqrt{13}}{3} - 1$}
\opt{C}{$\dfrac{\sqrt{17}}{3} - 2$}
\opt{D}{$\dfrac{2\sqrt{17}}{3} - 2$}
\end{qoptions}
\qn{10} $\displaystyle\int \dfrac{x+2}{(x^2+3x+3)\sqrt{x+1}}dx$ is equal to-
\begin{qoptions}
\opt{A}{$\dfrac{1}{\sqrt{3}}\cos^{-1}\left(\dfrac{x\sqrt{3}}{x+1}\right) + C$}
\opt{B}{$\dfrac{1}{\sqrt{3}}\tan^{-1}\left(\dfrac{x}{\sqrt{3(x+1)}}\right) + C$}
\opt{C}{$\dfrac{2}{\sqrt{3}}\tan^{-1}\left(\dfrac{x}{\sqrt{3(x+1)}}\right) + C$}
\opt{D}{$\dfrac{2}{\sqrt{3}}\cot^{-1}\left(\dfrac{x}{\sqrt{x+1}}\right) + C$}
\end{qoptions}
(where C is constant of integration)
\qn{11} The foot of perpendicular drawn from origin on line 3x + y = $\lambda$ is `P'. If line meets x-axis at `A' and y-axis at `B' then BP : PA is
\begin{qoptions}
\opt{A}{9 : 1}
\opt{B}{1 : 3}
\opt{C}{3 : 1}
\opt{D}{1 : 9}
\end{qoptions}
\qn{12} A normal to hyperbola 4x$^2$ -- 9y$^2$ = 36 meets the co-ordinateaxes x and y at A and B respectively. If parallelogram OAPB (where `O' is origin) is formed then locus of P is
\begin{qoptions}
\opt{A}{9x$^2$ + 4y$^2$ = 169}
\opt{B}{4x$^2$ + 9y$^2$ = 121}
\opt{C}{4x$^2$ -- 9y$^2$ = 121}
\opt{D}{9x$^2$ -- 4y$^2$ = 169}
\end{qoptions}
\qn{13} The number of four digit natural numbers which contains exactly two distinct digit
\begin{qoptions}
\opt{A}{567}
\opt{B}{576}
\opt{C}{657}
\opt{D}{675}
\end{qoptions}
\qn{14} Mean and variance of binomial probability distribution are 4 and $\dfrac{4}{3}$ respectively then probability of atleast two success is equal to
\begin{qoptions}
\opt{A}{$\dfrac{552}{729}$}
\opt{B}{$\dfrac{201}{243}$}
\opt{C}{$\dfrac{298}{343}$}
\opt{D}{$\dfrac{716}{729}$}
\end{qoptions}
\qn{15} The value of $\displaystyle\sum_{r=3}^{100}\dfrac{r}{(r-2)!+(r-1)!+r!}$ is-
\begin{qoptions}
\opt{A}{$\dfrac{100!+2}{2.100!}$}
\opt{B}{$\dfrac{100!-2}{2.100!}$}
\opt{C}{$\dfrac{100!-3}{2.200!}$}
\opt{D}{$\dfrac{100!+1}{100!.2}$}
\end{qoptions}
\qn{16} Area bounded by x$^2$ + y$^2$ $\leq$ 1 $\leq$ x + y is equal to-
\begin{qoptions}
\opt{A}{$\dfrac{\pi^2}{4}$}
\opt{B}{$\dfrac{\pi^2}{2}$}
\opt{C}{$\dfrac{\pi}{4}+\dfrac{1}{2}$}
\opt{D}{$\dfrac{\pi}{4}-\dfrac{1}{2}$}
\end{qoptions}
\qn{17} Let line `L' meets L$_1$ : $\dfrac{x}{1}=\dfrac{y}{2}=\dfrac{z}{3}$ and L$_2$ : $\dfrac{x-2}{2}=\dfrac{y-1}{4}=\dfrac{z-4}{5}$ orthogonally at point P and Q. If (PQ)$^2$ is D and DR's of line L are (a,b,c) where a,b,c $\in$ I, then $\dfrac{5D}{3}$ + $|$a$|$ + $|$b$|$ + $|$c$|$ can be
\begin{qoptions}
\opt{A}{4}
\opt{B}{6}
\opt{C}{8}
\opt{D}{10}
\end{qoptions}
\qn{18} Let f : [--1, 3] $\rightarrow$ [--8, 72] be defined as f(x) = 4x$^3$ -- 12x, then `f' is
\begin{qoptions}
\opt{A}{injective but not surjective}
\opt{B}{injective and surjective}
\opt{C}{neither injective nor surjective}
\opt{D}{surjective but not injective}
\end{qoptions}
\qn{19} If $|$z$_1$ -- 3 -- 4i$|$ = 1 and $|$z$_2$ -- 4 -- 3i$|$ = 1 then angle between locus of z$_1$ and z$_2$ is-
\begin{qoptions}
\opt{A}{$\dfrac{\pi}{3}$}
\opt{B}{$\dfrac{\pi}{6}$}
\opt{C}{$\dfrac{\pi}{2}$}
\opt{D}{$\dfrac{\pi}{4}$}
\end{qoptions}
\qn{20} A wire of length 2 unit is cut into two part which are bent to form square of side `x' and circle of radius `r'. If sum of areas of squares and circle so formed is minimum then x is equal to
\begin{qoptions}
\opt{A}{$\dfrac{2}{\pi+4}$}
\opt{B}{$\dfrac{\pi+4}{2}$}
\opt{C}{$\dfrac{4-\pi}{\pi}$}
\opt{D}{$\dfrac{2\pi}{\pi+8}$}
\end{qoptions}
\sect{SEC-2}
\qn{1} Let $\vec{a},\vec{b},\vec{c}$ are non collinear unit vectors such that $\vec{a}\times(\vec{b}\times\vec{c}) = \dfrac{\vec{b}+\vec{c}}{2}$. If angle between $\vec{a}$ and $\vec{b}$ is `$\alpha$' and angle between $\vec{a}$ and $\vec{c}$ is `$\beta$' then value of cos($\alpha-\beta$) -- cos($\alpha+\beta$) is -
\qn{2} If length of latus rectum of ellipse is 4 unit and the distance between a focus and its nearest vertex on major axis is $\dfrac{3}{2}$ unit then eccentricity of ellipse is
\qn{3} The length of subtangent to curve $\sqrt{x}+\sqrt{y}=3$ at point (4,1) is
\qn{4} If expression x + $\dfrac{1}{x^2}$, (x $>$ 0) attains its minimum value at x = $\alpha$, then $\alpha^3$ is-
\qn{5} A tower T$_1$ of height 60m is located exactly opposite to tower T$_2$ of height 80m on a straight road. From top of T$_1$, if angle of depression of foot of T$_2$ is twice the angle of elevation of top of T$_2$, then distance (in meter) between tower T$_1$ and tower T$_2$ is-
\qn{6} If $\displaystyle\lim_{x\to\sqrt{a}}\dfrac{b-\cos(x^2-a)}{(x^2-a)\sin(cx^2-a)} = $ L, is non-zero finite (a$>$0), then value of (L + b + c) is -
\qn{7} If $\alpha,\beta,\gamma$ are roots of equation x$^3$ -- 2x$^2$ -- 1 = 0 and a$_n$ = $\alpha^n+\beta^n+\gamma^n$, then value of $\dfrac{a_{11}-a_8}{a_{10}}$ is-
\qn{8} If $f(x) = \dfrac{\displaystyle\int_0^{x^2}\sin\sqrt{x}\,dx}{x^3}$, x $\geqslant$ 0 is continuous function at x = 0, then value of $f(0)$ is-
\qn{9} Let $\displaystyle\sum_{i=1}^{9}x_i^2 = 800$ and $\displaystyle\sum_{i=1}^{9}x_i = 80$ where x$_i$ is i$^{\text{th}}$ observation. If one observation that was 3, was found to be wrong and replaced by corrected value 4 then corrected variance is
\qn{10} Number of irrational terms in the expansion of $\left(3^{\frac{1}{2}}+5^{\frac{1}{3}}\right)^{60}$ is
\end{multicols}
\end{document}
